\subsection{control con convergencia en tiempo prescrito}

Considere el sistema no lineal

\begin{equation}
	\begin{aligned}
        \dot{x}_1&=x_2\\
	    \dot{x}_2&=f(x,t)+u\\
	    y&=x_1   
    \end{aligned}
\end{equation}

donde $|f(x,t)| \leq F$. Dado un tiempo $T>0$ se debe encontrar u tal que $y=r \forall t > T$ para la referencia $r$. Segun los errores de seguimiento

\begin{equation}
    \begin{aligned}
	    e_1&=r-y=r-x_1\\
	    e_2&=\dot{r}-x_2
    \end{aligned}
\end{equation}

derivando

\begin{equation}
    \begin{aligned}
	    \dot{e}_1&=\dot{r}-\dot{x}_1=\dot{r}-x_2=e_2\\
	    \dot{e}_2&=\ddot{r}-\dot{x}_2 = \ddot{r}-f(x,t)-u
    \end{aligned}
\end{equation}

de modo que

\begin{equation}
\begin{aligned}
	\dot{e}_1&=e_2 \\
	\dot{e}_2&=\ddot{r}-f(x,t)-u
\end{aligned}
\end{equation}

el problema de segumiento consiste en hacer que $e_1=e_2=0$ para todo $t\geq T$. Las condiciones iniciales del error son

\begin{equation}
\begin{aligned}
	e_1(0)=a\\
	e_2(0)=b
\end{aligned}
\end{equation}

Introducimos las trayectorias deseadas del error como funciones del tiempo, que deben cumplr

\begin{equation}
\begin{aligned}
	e_1^*(0)=a \quad e_1^*(T=0)\\
	e_2^*(0)=b \quad e_2^*(T=0)
\end{aligned}
\end{equation}

sea $\tau=\frac{t}{T}\epsilon[0,1]$  y la constante $c=bT+2a$, se propone el polinomio

\begin{equation}
\begin{aligned}
	e_1^*(t)=(1-\tau)^2(a+c\tau),\quad 0 \leq t \leq T
\end{aligned}
\end{equation}

con derivadas

\begin{equation}
\begin{aligned}
	e_2=& \dot{e}_1^* = b - \frac{2}{T} (3a+2bT)\tau + \frac{3C}{T}\tau^2\\
	\dot{e}_2=&\ddot{e}_1^*= - \frac{2}{T}^2(3a+2bT)+\frac{6C}{T^2}\tau
\end{aligned}
\end{equation}

Se deriva el error de seguimiento para las trayectorias deseadas del error

\begin{equation}
\begin{aligned}
	\epsilon_1=e_1^*-e_1\\
	\epsilon_2=e_2^*-e_2
\end{aligned}
\end{equation}

derivando

\begin{equation}
\begin{aligned}
	\dot{\epsilon}_1=&\dot{e}_1^*-\dot{e}_1=e_2-e_2=\epsilon_2\\
	\dot{\epsilon}_2=&\ddot{e}_1^*-\ddot{r}+f(x,t)+u
\end{aligned}
\end{equation}

con las condiciones iniciales

\begin{equation}
\begin{aligned}
	\epsilon_1(0)=0 \quad  y \quad \epsilon_2(0)=0
\end{aligned}
\end{equation}